\uuid{OJAZ}
\titre{Limite en un point}
\auteur{bodin ; hanani ; mzari}
\organisation{exo7}
\licence{CC BY-NC-SA 4.0}
\datecreate{2020-01-20}
\niveau{L1}
\module{Analyse}
\chapitre{Continuité, limite et étude de fonctions réelles}
\sousChapitre{Limite de fonctions}
\difficulte{1}
\theme{Limites des fonctions réelles, limite en un point}
\exo7qcmid{lille-1.limites.002}

\contenu{
\begin{question}
\qclassification{123.03}
\qtags{motcle=limite en un point}
Soit $f(x)= \frac{x^2-1}{2x^2-x-1}$. Quelles sont les assertions vraies ?
\begin{answers}
    \bad{$\lim_{x\to 1} f(x)=0$}

    \good{$\lim_{x\to 1} f(x)=\frac{2}{3}$}

    \bad{$\lim_{x\to -\frac{1}{2}} f(x) = +\infty$}

    \good{$\lim_{x\to (-\frac{1}{2})^+} f(x)=+\infty$}
\end{answers}
\begin{explanations}
$f(x)=\frac{(x-1)(x+1)}{(x-1)(2x+1)}= \frac{x+1}{2x+1}$.
\end{explanations}
\end{question}
}
